\scanpage{61}

\handwritten{\uline}{Conjecture} 1.

\handwritten{\tikz[baseline,overlay,x=1em,y=1em]{\draw[very thin,gray] (-1.55,-1.12) -- (-.72,-.51) -- (-.48,-.77) (-.76,-.44) -- (-.3,-.49);}}\handwritten{\uline}{Let} $q\handwritten{\leqslant}23$ be \handwritten{\uline}{prime},\handwritten{\uline}{and} $p$ \inserted{\handwritten{be}} a \handwritten{\uline}{prime other than} $q$,\handwritten{\uline}{and let $\handwritten{\underline}{\alpha}$\inserted{and $n$} be any positive integers} . \handwritten{\uline}{Let also}
\[
a(n)=c(q^{\handwritten{\underline}{\alpha}}n)/c(q^{\handwritten{\underline}{\alpha}})\ .
\]
\handwritten{\uline}{Then}
\[
\begin{aligned}
a(np)-a(n)a(p)+p^{-1}a(n/p)&\handwritten{\equiv}0 \qquad (\mathrm{mod}\ q^{\handwritten{\underline}{\alpha}})\qquad ,\\
a(nq)-a(n)a(q)&\handwritten{\equiv}0 \qquad (\mathrm{mod}\ q^{\handwritten{\underline}{\alpha}}).
\end{aligned}
\]
\margin{\handwritten{\#}}

In fact for $\handwritten{\underline}{q}\handwritten{\leqslant}5$ the modulus of the congruences can be improved to $q^{\handwritten{\underline}{\theta}}$\reversemarginpar\margin{\handwritten{theta}}\normalmarginpar with $\handwritten{\underline}{\theta}>\handwritten{\underline}{\alpha}$ . The reader familiar with the Ramanujan-type congruences(Lehner\handwritten{\uwave}{(5)},\handwritten{\uwave}{(6)},Atkin\handwritten{\uwave}{(1)})
\[
c(q^{\handwritten{\underline}{\alpha}}n)\handwritten{\equiv}0\quad(\mathrm{mod}\ q^{\handwritten{\underline}{\beta}})
\tag*{(3)}
\]
for $\handwritten{\underline}{q}=2,3,5,7,11$ and $\handwritten{\underline}{\beta}=3\handwritten{\underline}{\alpha}+8,2\handwritten{\underline}{\alpha}+3,\handwritten{\underline}{\alpha}+1,\handwritten{\underline}{\alpha},\handwritten{\underline}{\alpha}$\reversemarginpar\margin{\handwritten{beta}}\normalmarginpar respectively may query the definition of $\handwritten{\underline}{a(n)}$,but in fact for $\handwritten{\underline}{n}=1$ the result of (3) is best possible(i.\handwritten{e}. $\handwritten{\underline}{q^{\handwritten{\underline}{\beta}}\handwritten{\parallel}c(q^{\handwritten{\underline}{\alpha}})}$),and the division is therefore valid in any ring $\handwritten{\underline}{Z/q^{\handwritten{\underline}{\theta}}Z}$ . Next we state

\handwritten{\uline}{Conjecture} 2. \quad \handwritten{\uline}{Let} $\handwritten{\underline}{q=29\ \text{or}}$ $31$,\inserted{\handwritten{and}} $p$ \handwritten{\uline}{be a prime other than} $q$,\handwritten{\uline}{and $\handwritten{\underline}{\alpha}$ and $n$ be any positive integers}. \handwritten{\uline}{Then}
\[
c(q^{\handwritten{\underline}{\alpha}}n)/c(q^{\handwritten{\underline}{\alpha}})=\handwritten{\underline}{\handwritten{\lambda}}a(n)+\handwritten{\underline}{\handwritten{\bar{\lambda}}}\,\handwritten{\overline}{a(n)}\ ,
\]
\handwritten{\uline}{where}$\handwritten{\underline}{\handwritten{\lambda}},a(n)$ \handwritten{\uline}{and their conjugates}$\handwritten{\underline{\bar{\lambda}}},\handwritten{\overline}{a(n)}$ \handwritten{\uline}{lie in a quadratic extension} $E$ of $Z/q^{\handwritten{\underline}{\alpha}}Z$ \handwritten{\uline}{and}
\[
\begin{aligned}
a(np)-a(n)a(p)+p^{-1}a(n/p)&=0\ ,\\
a(nq)-a(n)a(q)&=0.
\end{aligned}
\]
\margin{\handwritten{\#}}
(We have written ``$=$'' rather than ``$\handwritten{\equiv}$'' in $E$ \ ). The situation for $\handwritten{\underline}{q}\handwritten{\geqslant}37$ appears to be that something like \deleted{c}\handwritten{C}onjecture 2 is valid; one finds that $\handwritten{\underline}{c(q^{\handwritten{\underline}{\alpha}}n)/c(q^{\handwritten{\underline}{\alpha}})}$ is a linear combination of ``congruence multiplicative'' coefficients lying in extensions of $\handwritten{\underline}{Z/q^{\handwritten{\underline}{\alpha}}Z}$ and the ``dimension'' of the``congruence space'' appears to be the integral part of $\handwritten{\underline}{q}/12$ (this is somewhat speculative) in the sense that the ``dimension'' is 1 in conjecture 1 and 2 in Conjecture 2.

We have so far considered the coefficient $\handwritten{\underline}{c(n)}$ on the full modular group. H\inserted{\handwritten{o}}wever the results quoted above give the prime $q$ a sp ecial role,\handwritten{\tikz[baseline,overlay,x=1em,y=1em]{\draw[very thin] (1.16,1.03) .. controls (1.29,1.31) and (1.46,1.16) .. (1.62,.69);}}both as to the definition of $\handwritten{\underline}{a(n)}$ and the multiplicative property; also $\handwritten{\sum\limits_{n=1}^{\infty}}c(q^{\handwritten{\underline}{\alpha}}n)\,x^{n}$ is (for all $\handwritten{\underline}{\alpha}$ ) a function on the subgroup $\handwritten{\underline}{\handwritten{\Gamma}}_{\handwritten{0}}(\handwritten{\underline}{q})$ defined by $\handwritten{\underline}{c\handwritten{\equiv}0\ (\mathrm{mod}\ q)}$ in

$z\ \rightarrow\ Vz=(az+b)/(cz+d)$. It therefore seems that these results are naturally associated with $\handwritten{\underline}{\handwritten{\Gamma}}_{\handwritten{0}}(q)$ and we now consider this

% archive sheet number: 60
